\chapter{Zeichensystem}
\label{ch:03_sign-system}

\begin{universityTask}
    Auch mit Bilden kann man lügen: Bildliche \q{Fake News} und Originalbild.

    \vspace{0.5cm}
    \begin{center}
        \begin{minipage}{0.49\textwidth}
            \centering
            \includegraphics[width=\textwidth]{t-00_sign-system-picture-fake.png}
        \end{minipage}
        \hfill
        \begin{minipage}{0.49\textwidth}
            \centering
            \includegraphics[width=\textwidth]{t-01_sign-system-picture-original.png}
        \end{minipage}
    \end{center}
    \vspace{0.5cm}

    \begin{enumerate}[label=(\alph*)]
        \item Welche Beziehungen bestehen zwischen den Begriffen Zeichensystem und Medium?
            \universityPoints{2}
        \item Wodurch unterscheiden sich sprachliche und bildliche Zeichensysteme vor allem i) syntaktisch und ii) semantisch voneinander?
            \universityPoints{4}
        \item Geben Sie ein Beispiel für einen markanten abbildungswertlichen Teil der Bildsyntax von Abbildung 2 und beschreiben Sie dieses Element anschließend möglichst genau eigenwertlich.
            \universityPoints{4}
        \item Geben Sie jeweils mindestens ein Beispiel dafür an, dass in Abbildung 2 indexikalische, ikonische und symbolische Zeichenaspekte wirksam sind (mit Begründung!).
            \universityPoints{6}
        \item Wie die beiden Bilder in Abbildung 2 ist auch der sprachliche Ausdruck \q{auf dem Dach eines Hochhauses in New York City} eine gut funktionierende Kontextbildung. Wodurch unterscheiden sich sprachliche und bildliche Kontextbildungen wesentlich?
            \universityPoints{4}
    \end{enumerate}
\end{universityTask}

\section{Zeichensystem und Medium}
\label{sec:03-01_sign-system-and-medium}

Bei Medien und Zeichen handelt es sich um Mittel zur Kommunikation. Als \textbf{Zeichensystem} bezeichnet man einen Komplext zusammengehöriger und alternativer Zeichenhandlungen, dessen Zusammengehörigkeit -- von weit zu eng -- sich orientiert an

\begin{itemize}
    \item der \textbf{Pragmatik}, bei der alle zusammenhängenden Handlungen betrachtet werden,
    \item der \textbf{Semantik}, welche lediglich die Bedeutung der Zeichen bestimmt, und
    \item der \textbf{Syntaktik}, die Beziehungen zwischen Zeichenträgern herstellt.
\end{itemize}

Zeichenintentionen werden mittels unterschiedlicher Zeichenträger innerhalb eines Mediums eingesetzt. \cite[S. 18 f.]{001_MKO02_media-as-means-of-communication}

Bei \textbf{Medien} handelt es sich um \qit{Mittel zur Übertragung und Verbreitung von Information durch Sprache, Schrift, Bild, Musik oder nonverbale […] Verständigungsweisen […]}, die klassifizierbar in die Klassen Eins bis Drei sind \cite[S. 7 f.]{001_MKO02_media-as-means-of-communication}. Dabei stammt der Begriff des Mediums aus dem Lateinischen, wo \qit{medium} in etwa so viel wie \qit{das Vermittelnde} bedeutet \cite[S. 6]{001_MKO02_media-as-means-of-communication}.

So wirken beide Begriffe zusammen: Medien vermitteln Zeichenhandlungen. Ein Zeichensystem wird innerhalb eines Mediums transportiert und erzeugt Bedeutung innerhalb dieses Mediums.

\section{Sprachliche und bildliche Zeichensysteme}
\label{sec:03-02_linguistic-and-pictorial-sign-systems}

\autoref{tab:03-01_syntax-and-semantics-of-linguistic-and-pictorial-sign-systems} auf Seite \pageref{tab:03-01_syntax-and-semantics-of-linguistic-and-pictorial-sign-systems} zeigt die wesentlichen Unterschiede zwischen sprachlichen und bildlichen Zeichensystemen in Bezug auf ihre Syntax und ihre Semantik.

\begin{table}[htp]
    \centering
    \begin{university-table}{width=\textwidth, colspec={l X[1,1] X[1,1]}}
        \textbf{} & \textbf{Sprachliche Zeichensysteme} & \textbf{Bildliche Zeichensysteme} \\
        \textbf{Syntax} & 
        Sprachliche Zeichensysteme bestehen aus einem endlichen Repertoire an Grundzeichen, die aneinandergereiht werden. Es wird zwischen gesprochener und geschriebener Alltagswortsprache unterschieden. Die gesprochene Sprache besteht aus Rede, welche Sätze enthält, die wiederum aus Wörtern bestehen, welche schließlich Laute enthalten. Die geschriebene Sprache lässt sich untergliedern in Texte aus Sätzen. Sätze enthalten Wörter, welche aus Buchstaben zusammengesetzt sind. \cite[S. 19 f.]{000_MKO01_communication-and-semiotics} &
        Bildliche Zeichensysteme bestehen ebenfalls aus elementaren Einheiten, allerdings sind diese anders als in der Sprache nicht hintereinander, sondern eher frei zusammengefügt. Sie können beispielsweise Linien oder Punkte sein. Es wird unterschieden zwischen einer eigenwertlichen Syntax, welche sich auf die Punkte, Linien, Flächen, Farben und Helligkeit eines Bildes konzentriert, und einer abbildungswertlichen Syntax, welche die Gesamtbetrachtung zusammengefügter Einzelelemente in den Fokus rückt. \cite[S. 24 f.]{000_MKO01_communication-and-semiotics} \linebreak
        Ein besonderes Merkmal bildlicher Zeichensysteme ist die sogenannte \qit{syntaktische Dichte}, die festlegt, dass Abweichungen zwischen Bildern im Prinzip beliebig klein sein können. \cite[S. 27]{000_MKO01_communication-and-semiotics} \\
        \textbf{Semantik} & 
        Die Semantik sprachlicher Zeichensysteme bestimmt die Bedeutung verwendeter Ausdrücke. Ein Defizit der Alltagssprachlichkeit ist, dass es ihr oft an Eindeutigkeit, semantischer Bestimmtheit und Schärfe beziehungsweise Klarheit mangelt. Daher bietet der Sprechakt als kommunikative Handlung (\gls{vgl} \autoref{sec:02-05_speech-act-and-language-game}) immer Raum für weitere Bedeutungsnuancen. \cite[S. 21 f.]{000_MKO01_communication-and-semiotics} &
        Die Bedeutung von Bildern kann meist über visuelle Ähnlichkeiten erkannt werden. Dabei finden ein Abgleich von Einzelgegenständen und die Bestimmung ihrer charakteristischen Attribute statt. Die syntaktischen Eigenwerte eines Bildes können dessen Semantik genauso prägen wie es die Gefühlswerte als Bedeutungsnuancen bei der Wortsprache tun. \cite[S. 26 f.]{000_MKO01_communication-and-semiotics} \\
    \end{university-table}
    \caption{Syntax und Semantik bei sprachlichen und bildlichen Zeichensystemen}
    \label{tab:03-01_syntax-and-semantics-of-linguistic-and-pictorial-sign-systems}
\end{table}

\section{Bildsyntax}
\label{sec:03-03_picture-syntax}

Anhand der Erkenntnisse zur Syntax bildlicher Zeichensysteme aus \autoref{tab:03-01_syntax-and-semantics-of-linguistic-and-pictorial-sign-systems} lässt sich nun die Bildsyntax der Abbildung aus der Aufgabenstellung genauer betrachten. Rein abbildungswertlich betrachtet stechen die Person und das Flugzeug heraus. Erstere, da sie schlicht einen besonders großen Teil des Bildes einnimmt, letzteres, da es als einziges Objekt nur in einer Version des Bildes zu sehen ist. Im Folgenden wird das Flugzeug eigenwertlich beschrieben.

Runde und glatte \textbf{Linien} grenzen das Objekt vom restlichen Teil des Bildes ab, teilweise grenzen sie auch einzelne Teile innerhalb des Objekts voneinander ab. Die Linien im oberen Bildteil sind unschärfer als die Linien im unteren Bildteil. Genau wie die Linien haben auch die \textbf{Flächen} des Objekts eine eher glatte Charakteristik. Es sind mehrere Teilflächen erkennbar, darunter eine größere ovale und zwei dazu orthogonal sehr längliche. Die \textbf{Farben} der Flächen sind Grau, Rot und Schwarz. Dabei ist ein Großteil des Objekts in hellem Grau gefärbt, ein farbiger roter Streifen befindet sich im unteren Teil des Objekts, neben ihm, symmetrisch angeordnet, befinden sich schwarze Kreise. Das Objekt grenzt sich durch eine auffallend starke \textbf{Helligkeit} von dem restlichen Bild ab, es ist eher hell und heller als der das Objekt umgebende Bildteil.

\section{Indexikalische, ikonische und symbolische Zeichenaspekte}
\label{sec:03-04_indexical-iconic-and-symbolic-aspects-of-signs}

\textbf{Indexikalische Aspekte} setzen Zeichen in ein Verhältnis zu Raum und Zeit. Mit Hilfe dieses Kontextes geben sie Auskunft über die raum-zeitliche Positionierung. \cite[S. 11 f.]{001_MKO02_media-as-means-of-communication} Geht man von der Integrität der Fotoaufnahme aus, so könnte -- ganz abstrakt gedacht -- jene selbst ein indexikalischer Aspekt sein, bezieht sie sich doch auf unmittelbar den Zeitpunkt ihrer eigenen Entstehung. Da das Bild, wie später noch betrachtet wird, den Zeitstempel \q{\texttt{09/11/01}} enthält, hebt dieser das Erstellungsdatum als Merkmal zusätzlich hervor.

\textbf{Ikonische Aspekte} stellen durch gemeinsame Eigenschaften oder bestimmte Ähnlichkeiten eine Verbindung vom Zeichen zum Objekt her. Sie ermöglichen es, dass ein Zeichen als das erkannt wird, was es darstellen soll. \cite[S. 11]{001_MKO02_media-as-means-of-communication} Im vorliegenden Bild gibt es viele Beispiele für ikonische Zeichenaspekte. Am einfachsten lässt sich sicherlich die Person als Person erkennen. Versierte Betrachter erkennen auch die Skyline der amerikanischen Stadt New York mit dem Hudson River. Aber auch das Flugzeug, bestehend aus den in \autoref{sec:03-03_picture-syntax} beschriebenen Linien und farbigen Flächen weißt durch die großflächigen Flügel sowie die runden und dunklen Triebwerke ausreichend Ähnlichkeit zum Gegenstand \q{Flugzeug} auf, um als solches erkannt zu werden.

\textbf{Symbolische Aspekte} beziehen sich mittels Konventionen auf bestimmte Sachverhalte. Diese Symbolik muss zunächst erlernt sein, um erkannt zu werden. \cite[S. 12]{001_MKO02_media-as-means-of-communication} Dass es sich bei den im unteren rechten Bildrand abgedruckten Zeichen um einen Datumsstempel handeln könnte, erkennt nur, wer diese Symbolik schon einmal erlernt hat. Dies ist ein besonders interessantes Beispiel für einen symbolischen Zeichenaspekt, da Datumsformate abhängig von Regionen und Ländern unterschiedlich interpretiert werden. So könnte \q{\texttt{09/11/01}}, angenommen es soll ein Datum darstellen, zum einen der elfte September oder der neunte November in amerikanischer oder in europäischer Schreibweise sein. Rein theoretisch müsste es nicht mal eine dieser beiden Optionen sein. Deutlich wird in jedem Fall, dass das Datum in Unkenntnis des korrekten Formats nicht interpretierbar wird.

\section{Kontextbildungen}
\label{sec:03-05_context-formations}

Bei Kontextbildung handelt es sich um eine Teilhandlung der sprachlichen Komunikation, die durch satzadverbiale Bestimmungen, beispielsweise \qit{auf dem Dach} oder \qit{in New York City}, aber auch durch Angaben zum Tempus der Verben realisiert wird \cite[S. 23 f.]{001_MKO02_media-as-means-of-communication}. Dabei wird der Kontext auch durch den Interaktionsbezug beeinflusst \cite[S. 22]{001_MKO02_media-as-means-of-communication}. Sprachlich ist Kontext direkt beschreibbar und theoretisch auch überprüfbar. Er wird Schritt für Schritt aufgebaut, während beim Leser ein mentales Bild des beschriebenen Sachverhalts konstruiert wird.

Bildzeichenhandlungen haben keinen direkten Sachbezug, mit den Propositionen fehlt auch die Einteilung in wahr oder falsch. Über Ikonen oder Symbole können allerdings durchaus Kontexte wahrgenommen werden \cite[S. 31 f.]{001_MKO02_media-as-means-of-communication}. Kontextbildung erfolgt bei Bildern also deutlich impliziter. Abgebildet wird eine Situation, in der das Abgebildete tatsächlich anwesend sein könnte -- im Bildraum wird dann die Situation evoziert. \cite[S. 30]{001_MKO02_media-as-means-of-communication} Bildlicher Kontext ist indirekt wahrnehmbar und kann zusätzlichen Kontext zur Überprüfung erfordern, er ist aber sofort simultan verfügbar.